\section{Verification Prompt 分析}

\subsection{Verifier 的角色与指令}

Verification 的 System Prompt 同样是结构化 Markdown，包含：

\begin{enumerate}
    \item \textbf{角色扮演}：你是卓越的数学家，IMO 级别严谨的审阅人
    \item \textbf{核心说明}：你\textbf{仅仅是 Verifier，不是 Solver}——不要尝试纠正错误或填充 gap，只需 step by step 验证
    \item \textbf{问题分类}：
    \begin{itemize}
        \item \textbf{Critical Error}：严重逻辑错误（如 $A > B, C > D$ 推出 $A-C > B-D$）
        \item \textbf{Gap}：跳步、逻辑断裂
    \end{itemize}
    \item \textbf{输出格式}：Summary（含 Verdict：valid/invalid）+ List of Findings（bug report）+ Detailed Verification Log
\end{enumerate}

\begin{warningbox}{Verifier 不应尝试解题}
Verification Prompt 明确要求 Verifier 不要尝试纠正错误或填充 gap，只做验证。这种\textbf{角色分离}是 Workflow 设计的关键——Generator 和 Verifier 各司其职。
\end{warningbox}

\subsection{本章小结}

Verification Prompt 通过严格的角色定义和结构化输出格式，确保 Verifier 只做验证、不做求解，其输出的 Verdict 和 Bug Report 驱动后续的 Correction 和循环终止判断。

